\scanpage{21}

l'ens des pts de $\operatorname{Spec} B$ où $B$ l'est est-il ouvert ?]. Il faut supposer que $A$ et $B$ sont des ``bons anneaux''. 

Tout d'abord, l'ens \unsure{précédent} est il \deleted{si $A$, $B$ sont locaux, $A\to B$ \unsure{local}} stable par localisation? \inserted{Cf IV 2.10.1}

% slanted
\margin{Ce qui est connu au cas où $A$ est \uline{intègre}, \unsure{primitivement|primitivem\textsuperscript{t}|prenons} \unsure{ou|que} (a) si $x\in U$, son $\unsure{Hc|H_0}$ primitive \unsure{au dessus du|au-dessus du} pt gén. de $y$ est $\in U$ \deleted{(b) Si $x\in U$, son \unsure{et constituant|il est} \unsure{gén.} Soit $x\in U$} (b) Soit $x$ le pt gén. de $y$, alors il existe un voisinage de $x$ ds $\bar{x}$ \deleted{qui} est dans $U$.}

% slanted
\diagram{\begin{tikzcd}[ampersand replacement=\&] x \arrow[d] \\ y \end{tikzcd}}

Cas d'extensions résiduelles finies. Aussi critère jacobien de Nagata et Th. de Nagata pour une algèbre locale noeth complète sur un \uline{corps}.

OK. en caractéristique $0$ pour les \uline{bons anneaux}, grâce à \uline{la résolution des singularités}.

% slanted
\margin{Pour (b), dans le cas d'un bon anneau, on est ramené au cas où $B$ est \uline{régulier}. Dans ce cas-ci on dispose de la résolution des singularités, et alors on le résout par l'\uline{affirmative}. \unsure{Concl.|Com.} : \unsure{le|b}, il se résout par l'affirmative pour les bons anneaux, \uline{en car. résiduelle nulle}, \deleted{\unsure{car cela}} \uline{pourvu qu'on sache} que le lieu où une $f$ sur $X$ est transversalement régulière est \uline{ouvert} (ou au moins constructible).}

4) $A\unsure{-|\to}B$ hom d'anneaux noeth. Lieu des pts de $\operatorname{Spec} B$ en lesquels $B$ est plat sur $A$ est-il ouvert? [\unsure{Ce est|Ceci est} \unsure{inconnu|inconnu.}: vérifier que pour $\forall\,\mathfrak{p}\in\operatorname{Spec} A$, et $\forall$ pt $\mathfrak{q}$ au dessus de $\mathfrak{p}$, maximal pour \unsure{ceci|cela}, posant $A'=A/\mathfrak{p}$, $B'=B/\mathfrak{q}$ \inserted{$A'\subset B'$, intègres} [dans $B'\otimes_{A'}K'$ \unsure{-|- -}] $\exists$ un ouvert non vide de $\operatorname{Spec} B'$ sur lequel $B'$ est plat sur $A'$.

% letterhead
M. DUFOUR\\
VARITYPISTE\\
131, RUE LECOURBE\\
PARIS XV\textsuperscript{e}\\
$\blacksquare$\\
C. C. P. Paris 14.143-34\\
R. M. 140.029
% archive sheet number: 20
